% =====================================================================
%  Chapitre 1 — Contexte général du projet
% =====================================================================
\chapter{Contexte général du projet}
\label{chap:contexte}

\section*{Introduction}
Ce premier chapitre pose le cadre du projet. Nous présentons d'abord
l'organisme d'accueil au sein duquel le stage a été réalisé, puis le
projet lui-même : sa problématique, une étude critique des solutions
existantes sur le marché et la solution que nous proposons. Nous concluons
par la présentation et la justification de la méthodologie de travail
adoptée.

\section{Présentation de l'organisme d'accueil}

\subsection{Le groupe Winvest Capital}
Le présent projet de fin d'études a été réalisé au sein de
\textbf{France Téléphone}, entité rattachée au groupe
\textbf{Winvest Capital}. Winvest Capital est un groupe
[À COMPLÉTER — secteur / activité principale du groupe, ex. investissement
et services numériques] regroupant plusieurs entités opérationnelles, dont
France Téléphone.

% [FIGURE : logo du groupe Winvest Capital]
\figplaceholder{Logo du groupe Winvest Capital}{fig:logo-winvest}{3cm}

\begin{description}[leftmargin=2.2cm,style=nextline,font=\bfseries]
  \item[Raison sociale] [À COMPLÉTER]
  \item[Date de création] [À COMPLÉTER]
  \item[Siège social] [À COMPLÉTER]
  \item[Secteur d'activité] [À COMPLÉTER]
  \item[Effectif] [À COMPLÉTER]
\end{description}

\subsection{France Téléphone}
\textbf{France Téléphone} est l'entité du groupe au sein de laquelle s'est
déroulé le stage. Son activité s'inscrit dans le domaine
[À COMPLÉTER — ex. de la distribution de forfaits et de services de
télécommunication], ce qui l'amène à produire de manière régulière des
supports de communication commerciale : offres promotionnelles, campagnes
e-mail, factures, contrats d'abonnement et messages à destination de sa
clientèle.

% [FIGURE : logo France Téléphone]
\figplaceholder{Logo de France Téléphone}{fig:logo-ft}{3cm}

Cette intensité de communication commerciale, sur des canaux multiples et
à un rythme soutenu, constitue précisément le terrain qui a fait émerger le
besoin à l'origine de ce projet.

\subsection{Cadre du stage}
Le stage s'est déroulé sur une durée de [À COMPLÉTER — ex. six mois], au
sein de l'équipe [À COMPLÉTER — ex. développement / systèmes
d'information]. La mission confiée consistait à concevoir et développer,
de bout en bout, une plateforme interne et commercialisable de création de
modèles de communication assistée par intelligence artificielle. Le travail
a couvert l'ensemble du cycle de vie logiciel : analyse des besoins,
conception, développement, tests et mise en production.

\section{Présentation du projet}

\subsection{Problématique}
La production de supports de communication professionnels se heurte à
plusieurs difficultés récurrentes que l'on peut synthétiser ainsi.

\begin{itemize}[leftmargin=1.4cm]
  \item \textbf{Une fabrication lente et technique.} La création d'un
  e-mail responsive de qualité suppose la maîtrise du HTML ou du MJML et
  une bonne connaissance des contraintes de compatibilité entre clients de
  messagerie. Les documents structurés — factures, contrats — imposent en
  outre une mise en forme rigoureuse et paginée.

  \item \textbf{Une dépendance à des compétences croisées.} Obtenir un
  résultat à la fois esthétique et techniquement correct requiert souvent
  la collaboration de profils créatifs (design) et de profils techniques
  (intégration), ce qui allonge les délais et alourdit le processus.

  \item \textbf{La multiplicité des canaux.} Une même campagne peut devoir
  être déclinée en e-mail, en SMS, en message RCS, voire accompagnée d'une
  facture ou d'un contrat. Or les outils du marché sont le plus souvent
  spécialisés sur un seul canal.

  \item \textbf{Le cloisonnement des données.} Dans un contexte
  multi-entreprises, chaque société doit disposer d'un espace de travail
  strictement isolé : ses modèles, ses utilisateurs et son identité de
  marque ne doivent jamais être visibles par une autre.

  \item \textbf{La rigidité des assistants existants.} Les générateurs
  automatiques disponibles produisent généralement du code HTML figé,
  difficile à reprendre et à modifier finement, plutôt que des éléments
  structurés directement éditables.
\end{itemize}

La problématique centrale du projet peut dès lors se formuler ainsi :
\emph{comment permettre à une équipe métier, sans compétence technique
particulière, de produire rapidement des supports de communication
multi-canaux, cohérents avec l'identité de son entreprise, tout en
garantissant l'isolation des données propre à un contexte
multi-entreprises et en tirant parti de l'intelligence artificielle sans
compromettre la souveraineté de ces données~?}

\subsection{Étude de l'existant}
Afin de situer notre solution, nous avons analysé plusieurs catégories
d'outils existants représentatifs du marché.

\begin{itemize}[leftmargin=1.4cm]
  \item \textbf{Canva} — Outil de conception graphique généraliste très
  répandu. Il excelle pour produire des visuels, mais n'est pas pensé pour
  l'e-mail responsive (pas de sortie MJML/HTML native robuste), ne couvre
  pas les documents transactionnels (facture, contrat) ni la logique
  multi-entreprises, et son assistance IA reste orientée image plutôt que
  génération de structures éditables par canal.

  \item \textbf{Mailchimp} — Plateforme d'e-mailing reconnue, dotée d'un
  éditeur par blocs et de fonctions marketing. Elle demeure toutefois
  centrée sur un seul canal (l'e-mail), pensée pour un compte unique plutôt
  que pour une architecture multi-locataire déléguée, et ne propose pas de
  reconstruction de campagne à partir d'une image.

  \item \textbf{Stripo et Beefree} — Éditeurs d'e-mails par glisser-déposer
  produisant du HTML/MJML de bonne qualité et exportables. Ils offrent une
  expérience visuelle soignée, mais ne fournissent pas d'assistant
  \emph{agentique} générant de vrais blocs éditables, ne couvrent pas
  l'ensemble des canaux (SMS, RCS, facture, contrat) et n'intègrent pas de
  fonction «~affiche vers e-mail~».

  \item \textbf{Éditeurs MJML bruts} — Des outils tels que l'éditeur
  officiel MJML s'adressent à des développeurs et supposent l'écriture
  directe de balises~; ils sont donc inaccessibles à un utilisateur métier.

  \item \textbf{Générateurs d'e-mails par IA} — Certains services génèrent
  un e-mail à partir d'une invite en langage naturel, mais restituent un
  bloc de HTML monolithique, difficile à retoucher, sans modèle de données
  structuré ni contrôle fin du rendu.
\end{itemize}

\subsection{Critique de l'existant}
De cette étude se dégage un constat clair : aucun des outils analysés ne
réunit simultanément l'ensemble des propriétés requises par notre contexte.
Le tableau~\ref{tab:comparatif-existant} synthétise cette comparaison.

\begin{table}[H]
\centering
\caption{Comparaison synthétique des solutions existantes}
\label{tab:comparatif-existant}
\footnotesize
\setlength{\tabcolsep}{4pt}
\begin{tabularx}{\linewidth}{Xccccc}
\toprule
\textbf{Critère} & \textbf{Canva} & \textbf{Mailchimp} & \textbf{Stripo /}
& \textbf{Éditeurs} & \textbf{Winaity} \\
 & & & \textbf{Beefree} & \textbf{MJML} & \\
\midrule
Éditeur visuel par blocs        & \checkmark & \checkmark & \checkmark & --         & \checkmark \\
Multi-canal (5 types)           & --         & --         & --         & --         & \checkmark \\
Multi-entreprises (tenants)     & --         & partiel    & --         & --         & \checkmark \\
IA agentique (blocs éditables)  & --         & partiel    & --         & --         & \checkmark \\
Affiche $\rightarrow$ e-mail    & --         & --         & --         & --         & \checkmark \\
LLM auto-hébergé (souveraineté) & --         & --         & --         & --         & \checkmark \\
Accessible sans compétence tech.& \checkmark & \checkmark & \checkmark & --         & \checkmark \\
\bottomrule
\end{tabularx}
\end{table}

Les principales lacunes identifiées sont les suivantes~:
\begin{itemize}[leftmargin=1.4cm]
  \item une \textbf{spécialisation mono-canal} qui oblige à multiplier les
  outils~;
  \item l'\textbf{absence d'un vrai modèle de blocs éditables} partagé
  entre l'édition manuelle et l'édition assistée par IA~;
  \item l'\textbf{absence d'isolation multi-entreprises déléguée} adaptée à
  un éditeur de logiciels souhaitant l'intégrer à ses propres outils~;
  \item la \textbf{dépendance à des API d'IA tierces}, incompatible avec les
  exigences de confidentialité des données d'une entreprise.
\end{itemize}

\subsection{Solution proposée}
Pour répondre à ces limites, nous proposons \textbf{Winaity Template
Builder} : une plateforme SaaS multi-entreprises de création de modèles de
communication, articulée autour de trois idées directrices.

\begin{enumerate}[leftmargin=1.4cm]
  \item \textbf{Un éditeur visuel multi-canal unique.} Cinq types de
  modèles — e-mail, facture, contrat, SMS et RCS — sont créés depuis une
  même plateforme, chacun disposant d'un éditeur dédié. L'éditeur d'e-mail,
  produit phare, repose sur un canevas de blocs manipulés par
  glisser-déposer et sur un modèle de données structuré.

  \item \textbf{Un assistant d'IA agentique produisant de vrais blocs.} Au
  lieu de générer du HTML jetable, l'assistant appelle des \emph{outils} de
  construction qui renvoient exactement le même format de blocs que
  l'éditeur manuel. Édition visuelle et édition par IA sont donc
  parfaitement interchangeables. L'assistant sait générer un e-mail complet
  à partir d'une consigne, le modifier en conversation, ou reconstruire une
  campagne à partir d'une \emph{affiche} lue par vision par ordinateur.

  \item \textbf{Une plateforme d'entreprise souveraine.} L'ensemble des
  données est cloisonné par entreprise (\emph{tenant}). Le modèle de langage
  qui alimente l'IA est \textbf{auto-hébergé}, ce qui garantit que les
  données ne quittent jamais l'infrastructure du groupe. La plateforme
  intègre par ailleurs la gestion des rôles, un système d'abonnements et
  une couche d'intégration permettant à des outils externes de l'embarquer.
\end{enumerate}

\section{Méthodologie de travail}

\subsection{Choix d'une méthodologie agile}
Le projet se caractérise par un périmètre fonctionnel large, des exigences
susceptibles d'évoluer au fil des retours de l'entreprise, et une forte
composante d'exploration technique (notamment sur le volet intelligence
artificielle). Ces caractéristiques rendent inadaptées les approches de
développement séquentielles, dites «~en cascade~», qui figent les besoins
en début de projet et n'autorisent que tardivement la validation par le
client. Nous avons donc opté pour une \textbf{approche agile}, qui privilégie
les livraisons incrémentales, les retours fréquents et l'adaptation
continue. Le tableau~\ref{tab:agile-vs-cascade} résume cette comparaison.

\begin{table}[H]
\centering
\caption{Approche agile versus approche traditionnelle}
\label{tab:agile-vs-cascade}
\small
\begin{tabularx}{\linewidth}{lXX}
\toprule
\textbf{Aspect} & \textbf{Approche agile} & \textbf{Approche en cascade} \\
\midrule
Cycle de développement & Itératif et incrémental & Séquentiel et linéaire \\
Gestion du changement & Adaptation continue & Faible tolérance au changement \\
Livraison & Incréments fréquents & Livraison en fin de projet \\
Implication du client & Retours réguliers & Implication tardive \\
Gestion du risque & Réduite par itérations & Concentrée en fin de cycle \\
Documentation & Progressive & Exhaustive en amont \\
\bottomrule
\end{tabularx}
\end{table}

\subsection{Choix de Scrum}
Parmi les méthodes agiles (Scrum, Extreme Programming, Kanban), nous avons
retenu \textbf{Scrum} pour la clarté de son cadre organisationnel et son
adéquation à un projet mené sur plusieurs mois. Scrum structure le travail
en itérations de durée fixe appelées \emph{sprints}, à l'issue desquelles
un incrément potentiellement livrable est produit. Il définit trois rôles :

\begin{itemize}[leftmargin=1.4cm]
  \item le \textbf{Product Owner}, qui porte la vision du produit et
  priorise le \emph{backlog}~;
  \item le \textbf{Scrum Master}, qui veille au bon déroulement du cadre et
  lève les obstacles~;
  \item l'\textbf{équipe de développement}, qui réalise l'incrément.
\end{itemize}

Dans le contexte particulier de ce stage, ces rôles ont été adaptés à une
configuration réduite : l'encadrant professionnel a joué le rôle de Product
Owner en exprimant et en priorisant les besoins, tandis que la réalisation
technique a été assurée dans le cadre du stage. Des points d'avancement
réguliers ont tenu lieu de cérémonies Scrum (planification, revue de
sprint, rétrospective).

% [FIGURE : schéma du processus Scrum — backlog, sprint, incrément]
\figplaceholder{Le processus Scrum}{fig:scrum}{5cm}

Le projet a été découpé en \textbf{quatre sprints}, chacun correspondant à
un ensemble cohérent de fonctionnalités~: (1) authentification, isolation
multi-entreprises et gestion des utilisateurs~; (2) éditeur de modèles
multi-canal~; (3) assistant d'intelligence artificielle agentique~; (4)
abonnements, intégrations et mise en production. Ce découpage est détaillé
au chapitre suivant.

\section*{Conclusion}
Ce chapitre a introduit l'organisme d'accueil, France Téléphone au sein du
groupe Winvest Capital, puis a posé la problématique du projet à travers
une étude critique des solutions existantes, faisant ressortir l'absence
d'un outil combinant approche multi-canal, isolation multi-entreprises,
intelligence artificielle agentique et souveraineté des données. Nous avons
enfin justifié le choix de la méthodologie Scrum. Le chapitre suivant
détaille l'analyse des besoins et la conception de la plateforme.
